\documentclass[11pt,a4paper]{article}
\usepackage[margin=24mm]{geometry}
\usepackage{amsmath,amssymb,amsthm,booktabs,microtype}
\usepackage[T1]{fontenc}
\usepackage{lmodern}
\usepackage[colorlinks=true,linkcolor=blue,urlcolor=blue,citecolor=blue]{hyperref}
\hypersetup{pdftitle={Subset Sums: optimized proof paths and a checked upper bound of 98},pdfauthor={Sum-of-Subset research project}}
\newcommand{\FS}{\operatorname{FS}}
\newcommand{\Z}{\mathbb Z}
\newcommand{\T}{\mathcal T}
\newtheorem{theorem}{Theorem}[section]
\newtheorem{lemma}[theorem]{Lemma}
\newtheorem{proposition}[theorem]{Proposition}
\title{Subset Sums in a Fourfold Sumset\\\large An optimized proof path from 99 to 98}
\author{Sum-of-Subset research project}
\date{16 October 2026}
\begin{document}
\maketitle
\begin{abstract}
We exclude cardinality 99 and continue the independently checked upper
bound to $N\le98$. The proof retains the residual loss after two empty
covering units, bounds core spans using small defect, and bounds the
inner double sumset by nested interval envelopes. A symbolic quadratic
checker replaces repeated checks at every integer side size: 5,443
interval certificates cover all 82,257 models, using 90 freshly generated
integer vectors. An independent execution mode expands every interval
and verifies each integer model. The resulting checked range is
$71\le N\le98$. The archived claim $N\le84$ remains stronger numerically;
this report advances the independently audited chain and simplifies its
finite certification, without claiming to surpass that archived record.
\end{abstract}
\tableofcontents

\section{Conclusion, dependencies, and the shorter proof path}
For a finite real set $A$, let
\[
 \FS(A)=\left\{\sum_{a\in X}a:X\subseteq A\right\},
 \qquad 4A=A+A+A+A.
\]
The empty sum is included and repeated summands in $4A$ are allowed.
Let $N$ denote the largest feasible cardinality.
\begin{theorem}\label{thm:main}
If $\FS(A)\subseteq4A$, then $|A|\le98$.
\end{theorem}
The known feasible set
$[-31,31]_{\Z}\cup\{\pm93,\pm155,\pm217,\pm713\}$ has 71 elements.
Thus the independently closed interval is now $71\le N\le98$.

Our foundation is the audited chain
$115\to103\to100\to99$ \cite{base103,base100,base99}.
We use its integralization, sign bootstrap, inverse theorem, propagation,
peeling budget, effective-hole lemma, and covering-unit lemmas.
The relevant lattice conditions are checked again at size 99.
The archived report \cite{archive} already contains the strategy for
this endpoint. We reconstruct its necessary inequalities, supply fresh
certificates, and reorganize the verification; we do not import its
endpoint certificate files. An arithmetic PASS is not a formalization
of the structural theorems: those remain mathematical dependencies.

The proof path has three stages. First the scalar budget eliminates high
holes and all configurations with no empty unit. Second a single
small-defect span lemma supplies the extra inequalities for both
remaining difficult regimes. Third, symbolic right-hand sides reduce
verification over integer parameter ranges to exact quadratic extrema.
The double-exception regime uses corrected rows at all retained centers,
so no enumeration of pairs of exceptional locations is needed.
The two discrete refinements are explicit; neither is hidden in a
numerical feasibility test.

\section{The size-99 reduction and scalar budget}
A hypothetical 99-element feasible set contains zero, since adjoining
zero otherwise gives a forbidden 100-element set. Apply the integralization
argument from the foundation, remove the four outer elements of each
sign to obtain $D$, and make two four-extreme peels to a core
$E=P\cup\{0\}\cup(-Q)$. With $T_r=r(r+1)/2$, put
\begin{gather*}
 d=|D|=91,\quad e=|E|=83,\quad a=|P|,\quad b=|Q|,\quad a+b=82,\\
 \{a,b\}=\{41-k,41+k\},\quad I=k^2,\quad
 \Lambda=1+T_a+T_b=1723+I,\quad\Omega=376-I.
\end{gather*}
The one-sided alternative is excluded by $1+T_{83}>24\cdot99+77$.
The inherited paired bootstrap uses
$q^2\le96e+1073-e^2$, $q\equiv e-1\pmod2$, and
$p=(e-1-q)/2$; it gives both side sizes at least 18.
For every oriented pair in this bootstrap range with $\Omega\ge0$, check
\begin{equation}\label{eq:gates}
 T_{\min(a,b)}-2-\Omega>0,\quad
 2(a+b)+4e+14-\Lambda+3\Omega<0,\quad
 \Lambda-(4d-10+\Omega)>0.
\end{equation}
There are 39 orientations; the respective extremal values are
196, $-85$, and 993. The common-lattice argument therefore applies.
After scaling, $D$ is integral. This does not yet give integrality of
all eight outer elements.

For a positive set $X$ of size $r$ define
\[
 \mu(X)=|\FS(X)|-1-T_r,\quad
 \delta_0(X)=|\{x\in X:2x\notin R(X\cup\{0\})\}|,
 \quad\rho(X)=\mu(X)+2-2\delta_0(X),
\]
where $R$ is the restricted two-sum set. Put $\rho=\rho(P)+\rho(Q)$ and
\[
 c=2|2D|+18d+85-|\FS(D)|,\qquad
 \kappa=|\FS(E)|-\Lambda-\mu(P)-\mu(Q).
\]
The inherited nonnegative slacks satisfy
\begin{equation}\label{eq:budget}
 \rho+2\theta+c+\T=376-I.
\end{equation}
Here $\T=\kappa+\ell_1+\ell_2+g$ consists of the core and peeling slacks.
Let $h_P,h_Q,z,w$ count the missing integers in the hulls of the
supports of $P,Q,E,D$, respectively. Then
\begin{gather}
 h_P,h_Q\le\kappa,\quad z=h_P+h_Q-\kappa,
 \quad0\le z\le\min(h_P,h_Q)\le\kappa,
 \quad h_X<|\FS(X)|-1,\label{eq:holes}\\
 w\le4z,\quad \T\ge\kappa+170-3z,
 \quad\T\ge\max\{57,z,170-2z\}.\label{eq:T}
\end{gather}
If $z>0$, neither side support is a progression. The foundation gives
$\mu(X)\ge r-3$ and $\rho(X)\ge r-14$ on either side, hence $\rho\ge54$.

\subsection{A small-defect bound for the high-hole stage}
\begin{lemma}\label{lem:small}
For a positive real set $X$ of size $r\ge15$, if
$0\le\rho(X)<5(r-14)$, then
\[
 \mu(X)\le\rho(X)+22+2\lfloor\rho(X)/5\rfloor
 \le\tfrac75\rho(X)+22.
\]
Consequently $\mu(X)\ge u$ implies
$\rho(X)\ge\min\{5(r-14),\lceil5(u-22)/7\rceil\}$.
\end{lemma}
\begin{proof}
For ordered prefixes $X_j$, use the inherited excess increments $y_j$
and recurrence $\rho(X_j)\ge\rho(X_{j-1})+y_j-2$.
If all $y_j\ge7$ for $15\le j\le r$, then
$\rho(X)\ge5(r-14)$, a contradiction.
Choose the last $j\ge15$ with $y_j\le6$. The small-increment inverse
lemma applies since $6\le j-5$; up to scaling, $X_j$ is an interval
$[1,j+q]\setminus H$ with $q\le6$ deleted positions.
Its diagonal-defect count is at most 12.
Indeed for $q\ge2$ all defects lie in $[1,q]$ or $[j,j+q]$.
There are at most $q-1$ in the first band: if all its positions were
present and defective, defectivity of 1 would force 2 absent.
There are at most $q+1$ in the second band. The one-deletion case has
at most two defects, and the interval has one.
Every later prefix adds at most one defect, while its increment is at
least 7. Nonnegativity of $\rho(X_j)$ gives
$r-j\le\lfloor\rho(X)/5\rfloor$. Thus
$\delta_0(X)\le12+\lfloor\rho(X)/5\rfloor$; substitute
$\mu=\rho-2+2\delta_0$. Separate $\rho\ge5(r-14)$ from its complement
to obtain the last assertion.
\end{proof}

The effective-hole lemma of \cite{base100} gives $\mu(X)\ge r+C+1$
when $2\beta<h_X<|\FS(X)|-1$ and the following exact tests hold:
\begin{gather*}
 K=\max\{k\ge3:k(k-3)/2\le C+1\},\quad D_*=C+K,\quad E_*=D_*+1,\\
 M=\max\left\{2\beta,1+\left\lfloor
 \frac{1+\sqrt{9+8(E_*+6\beta)}}2\right\rfloor\right\},\quad
 J=\max\{K+1,2\beta+5,M\},\\
 T_M-6\beta>M+1+E_*,\quad
 \beta(r-J+1)>D_*,\quad T_{J-1}-1>D_*.
\end{gather*}
As before, the finite search uses $20\le r\le64$, $7\le\beta\le20$,
and $0\le C\le600$. Every selected row is rechecked; optimality is
unnecessary. For each $z\ge15$, choose the greatest available $C$ with
$2\beta<z$ for each side, and set
\begin{align*}
 u_r&=r+C+1,\\
 v_r&=\max\{r-14,u_r+2-2r,
           \min(5(r-14),\lceil5(u_r-22)/7\rceil)\}.
\end{align*}
The second entry uses $\delta_0\le r$; the third uses Lemma~\ref{lem:small}.
Every actual state obeys
\begin{equation}\label{eq:cmax}
 c\le376-I-v_a-v_b-\max\{57,z,170-2z\}.
\end{equation}
The range $15\le z\le\Omega$ is exhaustive by \eqref{eq:holes}--\eqref{eq:budget}.

Write $R_p=\max P$, $R_q=\max Q$, $M_p=\max D$, $M_q=-\min D$,
and let $S_p,S_q$ be the positive and negative-magnitude masses of $D$.
Set $\varepsilon_p=\mu(P)+h_P$, $\varepsilon_q=\mu(Q)+h_Q$ and
\[
 F=S_p+S_q+c-2(M_p+M_q)-2(R_p+R_q).
\]
The inherited estimates give $F\le1764+4z$. If all eighteen covering
units are nonempty, after reflection the applicable lower alternative is
\begin{align*}
 F&\ge\tfrac12T_a+T_b+6a+8b+85
                  +\tfrac12\varepsilon_p+\varepsilon_q,\\
 \text{or }F&\ge\tfrac23T_a+T_b+6a+(22b+236)/3
                  +\tfrac23\varepsilon_p+\varepsilon_q.
\end{align*}
The verifier checks both alternatives in both orientations.
Use $\varepsilon_r\ge u_r+z$ at high holes, and use $u_r=0$ for $z=0$
or $u_r=r-3$ for $1\le z\le14$. After discarding negative budgets:
\begin{center}
\begin{tabular}{lrrr}
\toprule
Core holes&Surviving scalar states&Minimum active margin&Maximum high-hole $c$\\
\midrule
$0\le z\le14$&391&15&--\\
$z\ge15$&1217&$201/2$&73\\
\bottomrule
\end{tabular}
\end{center}
An empty unit costs $d=91$, so $c\le73$ forces all units active at high
holes. All high-hole states and all low-hole states with no empty unit
are excluded. There are 267 low-hole scalar states permitting empty
units. They are necessary relaxations, not asserted realizable sets.

\section{Classifying the remaining loss}
At zero holes the stronger budget is
\begin{equation}\label{eq:zero}
 \rho+\kappa+2\theta+c\le206-k^2.
\end{equation}
At positive low holes it gives $c\le152-k^2+2z\le180-k^2$.
There are at most two empty units, and two require $z=0$.
Put $\delta=c-91$ with one empty unit and $\delta=c-182$ with two.
Every retained center outside its target hull loses at least
$L=\min(a,b)+5=46-k$ points. The exhaustive regimes are
\begin{center}\small
\begin{tabular}{lrrrl}
\toprule
Regime&Range of $a$&Upper $\delta$&Upper $F$&Description\\
\midrule
Z&31--51&$115-k^2$&1764&One empty, no outside center\\
P&32--50&$89-k^2$&1820&Positive holes, no outside center\\
ZE&33--49&$115-k^2$&1764&One empty, one outside center\\
PE&34--48&$89-k^2$&1820&Positive holes, one outside center\\
ZZ&36--46&$115-k^2$&1764&One empty, two outside centers\\
EE&37--45&$24-k^2$&1764&Two empty, no outside center\\
\bottomrule
\end{tabular}
\end{center}
The ranges follow by imposing $\delta\ge0,L,2L$, respectively.
The identities
\[
 3L-(115-k^2)=(k-3/2)^2+83/4>0,\qquad
 2L-(89-k^2)=(k-1)^2+2>0
\]
exclude further exceptions. In EE, $24-k^2<L$.
With one empty unit, the inherited outer-integrality lemma makes every
outer gap at least one. In EE we use only weak outer order.

\subsection{One span lemma for both difficult regimes}
\begin{lemma}\label{lem:span}
Let $X\subset\Z_{>0}$ have size $r\ge16$ and full interval support.
If $0\le\rho(X)\le r-5$, then $\max X\le r+\rho(X)+1$.
\end{lemma}
\begin{proof}
Writing $x_1<\cdots<x_r$, fullness forces $x_1=1$ and
$x_i\le1+\sum_{j<i}x_j$; every prefix has full interval support.
Set $q=x_r-r\ge0$. The last excess increment is $y_r=q$.
If $\nu$ old diagonal defects disappear at the last step, the exact
recurrence is
\[
 \rho(X)=\rho(X_{r-1})+q-2+2\nu.
\]
Thus $q\le\rho(X)+2$. Equality would force
$\rho(X_{r-1})=\nu=0$ and $2\le q\le r-3$.
The zero-defect classification proved in \cite{base99} gives
$X_{r-1}=[1,r-1]$ or $[1,r]\setminus\{r-2\}$.
In the first case its defective maximum $r-1$ acquires the distinct-sum
representation $2(r-1)=x_r+(r-q-2)$.
In the second case $r,r-1$ are defective. For $q\ne2$ use
$2r=x_r+(r-q)$; for $q=2$ use $2(r-1)=x_r+(r-4)$.
The second summands are present in the stated prefixes, so in each case
$\nu>0$, a contradiction. Integrality gives the claimed sharper bound.
\end{proof}

In ZZ, \eqref{eq:zero} and $\delta\ge2L$ imply
$\rho+\kappa\le23+2k-k^2$, $k\le5$.
In EE they imply $\rho+\kappa+\delta\le24-k^2$, $k\le4$.
Both bounds on $\rho$ are strictly below $\min(a,b)-14$; the inherited
nonprogression penalty therefore forces both side supports to be
arithmetic progressions. A support progression of step at least two and
cardinality $s$ has at least $s-1$ hull holes. But
$h_X\le\kappa\le24<T_{\min(a,b)}\le s-1$.
Hence both supports have step one, and $h_P=h_Q=\kappa=0$.
The individual defect bounds meet Lemma~\ref{lem:span}; adding its two
conclusions to the budget yields
\begin{equation}\label{eq:span}
 R_p+R_q+\delta\le
 \begin{cases}199-k^2,&\text{ZZ},\\106-k^2,&\text{EE}.\end{cases}
\end{equation}

In ZZ each retained center loses at most
$\delta-L\le69+k-k^2\le69$. Its contribution to the hull is at least
$91-69=22>4$; at an outside endpoint it therefore has a contributing
core point. Counting the at most $R_p-a$ missing positive core positions
and $R_q-b$ missing negative core positions gives, for every positive
center $v$, the simultaneous relaxed rows
\begin{equation}\label{eq:ZZ}
 -v-l_v-R_p\le-87,\qquad v-u_v-S_p-R_q\le-86.
\end{equation}
Reflect for negative centers. This uses the actual spans, without
assuming the zero-defect bounds $R_p\le a+1$, $R_q\le b+1$.

\section{Mass and double-sumset information}
Distinct positive integers in the core give
$\sum P\ge T_{a-1}+R_p$. The next three points are at least
$R_p+1,R_p+2,R_p+3$ and the fourth is $M_p$. Consequently
\begin{equation}\label{eq:mass}
 S_p\ge T_{a-1}+4R_p+6+M_p,
 \qquad S_q\ge T_{b-1}+4R_q+6+M_q.
\end{equation}
These are imposed alongside the usual bounds
$S_p\ge T_a+\varepsilon_p+3R_p+6+M_p$ and its reflection.
The two lower bounds are not added together.

If $M_q-R_q=4$, the four negative outliers inside $D$ are consecutive.
Their ten unordered repeated pairs have only seven distinct sums.
All pairs except the same-sign outlier pairs lie in
$[-M_q-R_q,M_p+R_p]$. Therefore
\begin{equation}\label{eq:gap4}
 |2D|\le M_p+M_q+R_p+R_q+18,
 \qquad F\le1758\quad(z=0).
\end{equation}
The ordinary count uses $+21$ and gives $F\le1764$.
One ZE case will split the integral gap into 4 and at least 5.

In EE expose the four inner extreme magnitudes on each side:
\[
 x_0=R_p<x_1<x_2<x_3<x_4=M_p,\qquad
 y_0=R_q<y_1<y_2<y_3<y_4=M_q.
\]
All successive gaps are at least one, and
\begin{equation}\label{eq:innermass}
 S_p\ge T_{a-1}+\sum_{i=0}^4x_i,
 \qquad S_q\ge T_{b-1}+\sum_{i=0}^4y_i.
\end{equation}
\begin{lemma}[Nested interval envelopes]\label{lem:envelope}
For $0\le j\le3$, put $W_j=x_j+y_j$, $t_j=2(4-j)$. Then
\begin{equation}\label{eq:envelope}
 |2D|\le2W_j+1+\binom{t_j+1}{2}
 +\sum_{i=j+1}^4\!\left[
 \min(W_j+1,x_i-x_{i-1})+\min(W_j+1,y_i-y_{i-1})\right].
\end{equation}
\end{lemma}
\begin{proof}
Replace the core together with the first $j$ inner points on each side
by its integer hull $[-y_j,x_j]$. Its double sum has $2W_j+1$ integers.
Add the positive translates $x_i+[-y_j,x_j]$ in order. Their union with
the first interval enlarges at step $i$ by at most
$\min(W_j+1,x_i-x_{i-1})$: the preceding upper endpoint is $2x_j$
initially and $x_{i-1}+x_j$ subsequently. The negative translates give
the reflected bound. Adding the two enlargement bounds remains valid
even when the two sides overlap. Only pairs entirely among the $t_j$
external points remain; with repetition there are at most
$\binom{t_j+1}{2}$ of them. Every pair in $2D$ has now been covered.
\end{proof}
For each minimum introduce a nonnegative variable $g$ bounded above by
both arguments. An actual set realizes these rows by choosing the
minimum, so this is a valid relaxation. Set $K=|2D|$.
Since $w=0$ in EE, the definition of $c=182+\delta$ gives
\begin{equation}\label{eq:cardinality}
 S_p+S_q+\delta-2K=18\cdot91+84-182=1540.
\end{equation}
The model needs only its upper inequality.

\section{Complete specification of the finite models}
\subsection{Centers, paired units, and modes}
Let the outer magnitudes be $p_1>p_2>p_3>p_4>M_p$ and their $q$
reflection. Define
\[
 \lambda_p=p_1-p_2-p_3-p_4,\quad
 \alpha_p=p_2-p_3-p_4,\quad\beta_p=2p_2-p_1-p_3-p_4.
\]
The ten positive centers, indexed from zero, are
\[
 v^p=(\lambda_p+p_1,\ldots,\lambda_p+p_4,
       \lambda_p-q_1,\ldots,\lambda_p-q_4,\alpha_p,\beta_p).
\]
Interchange $p,q$ for $v^q$. Slots \texttt{p0},\ldots,\texttt{q9}
are center labels, distinct from the magnitude variables $p_i,q_i$.
Positive centers translate $D$ into $[1,S_p]$; negative centers
translate $-D$ into $[1,S_q]$.

There are sixteen single units and two paired units. Pair 1 uses
\texttt{p4,q4}. In branch A, $\beta_p+\beta_q\le0$, pair 2 uses
\texttt{p9,q9}; in branch B, $\beta_p+\beta_q>0$, it uses
\texttt{p5,q5}. Remove the chosen one or two missing units and retain
all other singles. If pair 1 is present, reflect so that its contributor
is \texttt{p4}. If pair 2 is present, choose its $p$ and $q$ contributors
in separate modes, except in branch A when \texttt{p5} or \texttt{q5}
is missing: then omit individual constraints on pair 2, allowing splitting.
No mode is duplicated for this dropped pair.

The foundation proves this choice exhaustive: pair 1 has at most one
contributor; so does pair 2 in B. In A the same follows when both
\texttt{p5,q5} contribute from
$\beta_p+\beta_q=-(v^p_5+v^q_5)-2(p_3+p_4+q_3+q_4)$.
There are 68 one-empty modes (33 A, 35 B) and 547 two-empty modes
(258 A, 289 B). In the one-empty modes, 66 retain 17 centers and two
retain 16. ZE and PE designate each retained center and each endpoint,
giving $2(66\cdot17+2\cdot16)=2308$ choices per side size.

\subsection{Rows shared by the regimes}
All variables are nonnegative real numbers. The base variables are
\[
 p_1,\ldots,p_4,q_1,\ldots,q_4,M_p,M_q,R_p,R_q,S_p,S_q,\delta
\]
and $(l_v,u_v)$ for each retained center. Rows have orientation $Ax\le b$.
Include the reflected counterpart of every positive-side row below:
\begin{align*}
 D_i^p &:p_{i+1}-p_i\le-1\quad(1\le i\le3),\qquad
 D_4^p :M_p-p_4\le-1,\\
 E_p &:R_p-M_p\le-4,\qquad G_p :-R_p\le-a,\\
 H_p &:2M_p+6R_p-2S_p\le-2T_a-12-2\eta_p,\\
 \widehat H_p &:M_p+4R_p-S_p\le-T_{a-1}-6.
\end{align*}
Here $(\eta_p,\eta_q)=(a-2,b-2)$ in P and PE, and $(0,0)$ otherwise.
In EE the right sides of all eight outer $D$ rows are zero.
The branch row is $\beta_p+\beta_q\le0$ in A, and
$-\beta_p-\beta_q\le-1$ in B; the B right side is zero in EE.
Add the regime's budget $\delta\le\delta_{\max}$ and
\begin{equation}\label{eq:loss}
 \sum_{v\text{ retained}}(l_v+u_v)\le\delta.
\end{equation}
For a positive in-hull center, the ordinary quantile rows are
\[
 -v-l_v\le-b-5,\qquad v-u_v-S_p\le-a-4.
\]
Reflect to obtain the negative-center rows.
For the designated outside endpoint in ZE or PE replace only that row:
\begin{center}\small
\begin{tabular}{lll}
\toprule
Endpoint&Replacement&Extra loss row\\
\midrule
Positive lower&$-v-l_v-M_p\le-91$&$-l_v\le-b-5$\\
Positive upper&$v-u_v-S_p-M_q\le-90$&$-u_v\le-a-5$\\
Negative lower&$-v-l_v-M_q\le-91$&$-l_v\le-a-5$\\
Negative upper&$v-u_v-S_q-M_p\le-90$&$-u_v\le-b-5$\\
\bottomrule
\end{tabular}
\end{center}
The opposite endpoint retains its ordinary row. No actual outside-location
condition is imposed, so the systems remain necessary relaxations.
In ZZ replace every quantile by \eqref{eq:ZZ} or its reflection, and add
$-\delta\le-2(46-k)$ and \eqref{eq:span}. EE retains ordinary quantiles
and adds its version of \eqref{eq:span}. Its residual budget is less than
$L$, so these centers are in-hull and contribute. A contributing center
is integral because $D$ and its target support are integral; outer
integrality is not needed.

Finally set
$f=S_p+S_q+\delta-2M_p-2M_q-2R_p-2R_q$ and include
$f\le F_{\max}-91$ for one-empty regimes or
$f\le F_{\max}-182$ for EE. Usually $F_{\max}$ is in the regime table;
the special gap-four branch uses 1758.

\subsection{EE envelope rows and two discrete refinements}
Only EE adds $x_1,x_2,x_3,y_1,y_2,y_3,K$ and
$g^p_{ji},g^q_{ji}$, $0\le j\le3$, $j<i\le4$.
With $x_0=R_p$, $x_4=M_p$ and the reflected $y$ notation, include
\begin{gather*}
 x_{i-1}-x_i\le-1\ (1\le i\le4),\qquad
 \sum_{i=0}^4x_i-S_p\le-T_{a-1},\\
 g^p_{ji}-x_i+x_{i-1}\le0,\qquad g^p_{ji}-W_j\le1,
\end{gather*}
and their reflections. Include the four envelope rows
\[
 K-2W_j-\sum_{i=j+1}^4(g^p_{ji}+g^q_{ji})
 \le1+\binom{t_j+1}{2},\quad0\le j\le3,
\]
and $S_p+S_q+\delta-2K\le1540$.
Two additional exhaustive splits finish the specification:
\begin{enumerate}
\item ZE, $a=47$, branch A, missing \texttt{q5}, exceptional lower
endpoint of \texttt{p4}: replace the single case by the two integral
possibilities $M_q-R_q\le4$ and $R_q-M_q\le-5$.
The ordinary row already gives $M_q-R_q\ge4$.
The first possibility uses $F\le1758$ by \eqref{eq:gap4}; the second
keeps $F\le1764$. Equality at 1764 is not treated as a contradiction.
\item EE, branch A, missing \texttt{p0,q5}: the dropped second pair has
combined contribution at least $91-\delta>0$.
At least one of its individual slots contributes. Its center is integral;
its largest translated point gives either
$\beta_p+M_p\ge1$ or $\beta_q+M_q\ge1$.
Replace each of the nine side sizes by two models, with the corresponding
row $-\beta_p-M_p\le-1$ or $-\beta_q-M_q\le-1$.
Both may hold; the disjunction is exhaustive without requiring a unique
contributor or integral outer elements.
\end{enumerate}
The complete counts are
\begin{center}
\begin{tabular}{lrr}
\toprule
Regime&Enumeration&Models\\
\midrule
Z&$21\cdot68$&1428\\
P&$19\cdot68$&1292\\
ZE&$17\cdot2308+1$&39237\\
PE&$15\cdot2308$&34620\\
ZZ&$11\cdot68$&748\\
EE&$9\cdot547+9$&4932\\
\midrule
Total&&82257\\
\bottomrule
\end{tabular}
\end{center}

\section{Symbolic interval verification and reproducibility}
\subsection{Why interval checking is exact}
Within a fixed structural group the matrix $A$ is constant and every
right-hand side is a rational polynomial of degree at most two in $a$:
$b=82-a$, $k^2=(41-a)^2$, and the triangular numbers are quadratic.
The only absolute value appears in the ZZ row $-\delta\le-2(46-k)$.
Split each ZZ group into $a\le41$ and $a\ge42$, so $k=41-a$ or $a-41$
is linear. The special ZE point is isolated at $a=47$, and the two EE
pair choices are separate groups. This produces 5,439 structural groups.

For integer weights $w_i>0$, verify $wA\ge0$ coordinate by coordinate.
Then form the polynomial $p(a)=wb(a)$ \emph{symbolically}, using exact
rational coefficient arithmetic. If $p(a)<0$ throughout the assigned
integer interval, no $x\ge0$ can satisfy $Ax\le b(a)$.
For $p(a)=c_0+c_1a+c_2a^2$, its integer maximum on $[L,U]$ occurs at
an endpoint when $c_2\ge0$; when $c_2<0$, also check the floor and ceiling
of $-c_1/(2c_2)$ when they lie in the interval.
Thus at most four exact evaluations suffice. There is no interpolation
assumption: the row formulas themselves are evaluated in a polynomial
class which rejects degree greater than two.

The certificate file records a pool of 90 integer vectors and a partition
of each group into closed integer intervals. The verifier reconstructs
the group universe from the mode rules, requires exact key equality,
checks contiguous coverage with no overlap or gap, and rejects unused
vectors. It checks 5,443 intervals instead of repeating matrix
combinations for all 82,257 integers, a reduction of about 93.4\%.
The stored result records every combined polynomial and its worst
integer parameter. The normalization of a contradiction is
$-\max p(a)/\sum_iw_i$; only strict positivity is accepted.
The minimum over all delivered intervals is $12/271$.

An optional \texttt{--exhaustive} mode independently builds the ordinary
integer model at every parameter, compares every matrix row and right-hand
side against the symbolic model, checks the integer combination, and
requires the same minimum margin as the quadratic extremum calculation.
This full cross-check was run for the delivered certificate set.

\subsection{Files and verification commands}
The report source and PDF are
\path{reports/upper98_from99_20261016.tex} and its PDF companion.
The standard-library checker is \texttt{verification/upper98.py}; the proof
data are \path{data/upper98_certificates.json}. The optional generator
\path{experiments/generate_upper98_certificates.py} uses SciPy only to
propose weights. Integer reconstruction and an exact check are required
before any proposed vector is accepted. Global certificate reuse required
91 optimizer calls to produce the delivered 90-vector pool.

From the repository root, using Python 3.10 or newer:
\begin{verbatim}
python -I -S -B verification/upper98.py --output .work/upper98.json
python -I -S -B verification/upper98.py --exhaustive
python -I -S -B verification/upper99.py
\end{verbatim}
On Windows, \texttt{verify\_upper98\_windows.bat} runs all four continuation
stages in order. Verification needs no optimizer, third-party Python
package, or network. The result, baseline rerun, generation log, and run
receipt are under \texttt{results/2026-10-16/upper98/}.

Ten negative controls reject negative weights, a noncontradictory vector,
an unknown row, an interior-vertex failure missed by endpoint-only checks,
a polynomial of excessive degree, a missing group, a missing integer,
an omitted gap-four branch, an omitted split-pair branch, and an attempt
to certify the unproved next cardinality. Checks use explicit exceptions,
not removable Python assertions. Normal and \texttt{-O} outputs agree
byte for byte; exhaustive and interval outputs agree apart from their mode
label. The preceding upper-99 finite result is reproduced unchanged.

\begin{proof}[Proof of Theorem~\ref{thm:main}]
The foundation excludes all sizes above 99. A feasible set of size 99
passes the sign and lattice reduction, hence satisfies the scalar budget.
The high-hole and all-active comparisons eliminate those cases.
Every remaining set maps, after the justified reflection and disjunctions,
to a model in one of the six complete regimes. All such models have a
strict nonnegative integer-row contradiction, certified throughout their
parameter intervals and also checked pointwise in the delivered run.
No feasible 99-element set exists.
\end{proof}

The next task is to exclude size 98; none of these size-specific budgets,
case counts, or certificates establishes that exclusion. The verifier
rejects it explicitly. Nor does this continuation complete the independent
audit of the stronger archived upper bound 84.

\begin{thebibliography}{9}
\bibitem{base103} Sum-of-Subset project,
\emph{A checked continuation from 115 to 103}, 14 October 2026.
Repository: \texttt{reports/upper103\_from115\_20261014.tex}.
\bibitem{base100} Sum-of-Subset project,
\emph{A checked continuation from 103 to 100}, 16 October 2026.
Repository: \texttt{reports/upper100\_from103\_20261016.tex}.
\bibitem{base99} Sum-of-Subset project,
\emph{A checked continuation from 100 to 99}, 16 October 2026.
Repository: \texttt{reports/upper99\_from100\_20261016.tex}.
\bibitem{archive} Sum-of-Subset project,
\emph{Complete upper-bound 86 report}, 5 October 2026,
section ``Retaining residual loss: a general upper bound of 98'' and
its certificate appendix.
Repository: \texttt{reports/upper86\_complete\_20261005.tex}.
\end{thebibliography}
\end{document}
